\section{Finite entropy, conditioning, and passage to profiles}
\label{sec:entropy}

The geometric arguments will use conditional point and parameter
populations whose masses must remain controlled after labels are fixed.
We first develop the finite probability estimates that justify those
operations, and then explain their passage to limiting entropy profiles.
All probability spaces in the finite statements are finite, with
cardinalities allowed to depend on $\eps$. Put $\ell=\log(1/\eps)$ and
write $\Entn,\Mutn$ for natural-log entropy and mutual information,
with $\Ent=\Entn/\ell$ and $\Mut=\Mutn/\ell$.
The exact particle index $T$ may have an arbitrarily large alphabet;
no bound on its entropy is required. Every comparison measure is
specified explicitly and need not be a marginal of the law being estimated.

\subsection{Entropy estimates with unrestricted parent alphabets}

Our conventions are the usual finite entropy and relative entropy
\citep{Shannon1948,KullbackLeibler1951}. Conditional copies and entropy
submodularity also feature in Tao's entropy calculus
\citep{Tao2010Entropy}; the comparison and
retention bounds required here are proved below.

We use the chain rule, nonnegativity of conditional mutual information,
and $\Entn(X)\le\log|\mathcal X|$ for an $\mathcal X$-valued observation.
Here are the particular consequences, including quantitative restriction
estimates, that we will need.

\begin{lemma}[Determination, masses, and restriction]
\label{lem:entropy-restriction}
Let $A,B,D,T$ be finite observations.
Then
\begin{align}
 \Mutn(T;B)&\le\Mutn(T;A)+\Entn(B\mid A),\label{eq:entropy-determination}\\
 |\Entn(B\mid T)-\Entn(A\mid T)|
 &\le\Entn(B\mid A)+\Entn(A\mid B).\label{eq:entropy-change}
\end{align}
If, outside an event of probability $\zeta$, $B$ has at most $N$ possible
values given $A$, and $|\mathcal B|\le M$, then
\begin{equation}
 \Entn(B\mid A)\le\log2+\log N+\zeta\log M.
 \label{eq:entropy-ambiguity}
\end{equation}
For nonnegative numbers $m_b$ with $m_b>0$ on the support of $B$,
\begin{equation}
 \Entn(B)+\mathbb E\log m_B\le\log\sum_bm_b.
 \label{eq:entropy-masses}
\end{equation}
Finally, let $\mathsf P,\mathsf Q$ be laws on the same space with
$\mathsf Q\le K\mathsf P$, $K\ge1$.  For arbitrary $B,D$,
\begin{align}
 \Entn^{\mathsf Q}(D\mid B)
 &=\mathbb E_{\mathsf Q}[-\log\mathsf P(D\mid B)]
   -\mathbb E_{\mathsf Q_B}
     \KL(\mathsf Q_{D\mid B}\Vert\mathsf P_{D\mid B}),\label{eq:conditional-kl}\\
 \mathbb E_{\mathsf Q_B}
     \KL(\mathsf Q_{D\mid B}\Vert\mathsf P_{D\mid B})
 &\le\log K.\label{eq:conditional-kl-bound}
\end{align}
In particular, restriction to an event of $\mathsf P$-mass $m>0$ costs
at most $\log(1/m)$ in \eqref{eq:conditional-kl-bound}.
None of these estimates depends on the cardinality of $B$ or $T$.
\end{lemma}

\begin{proof}
The chain rule gives
$\Mutn(T;B)\le\Mutn(T;A,B)
=\Mutn(T;A)+\Mutn(T;B\mid A)$, proving
\eqref{eq:entropy-determination}.  Apply the chain rule to
$\Entn(A,B\mid T)$ in both orders to obtain
\eqref{eq:entropy-change}.  Adjoining the indicator of the exceptional
event costs at most $\log2$; conditioning on its two values gives
\eqref{eq:entropy-ambiguity}.  For \eqref{eq:entropy-masses}, normalize
$m$ to a probability law and use nonnegativity of its relative entropy
from the law of $B$.  Expanding the conditional relative entropy gives
\eqref{eq:conditional-kl}.  Its expectation is at most
$\KL(\mathsf Q_{B,D}\Vert\mathsf P_{B,D})$, which by marginalization
is at most $\KL(\mathsf Q\Vert\mathsf P)\le\log K$.
\end{proof}

We call a determination \emph{of zero cost} when the right side of
\eqref{eq:entropy-ambiguity}, divided by $\ell$, tends to zero.
This includes bounded ambiguity between comparable translated grids,
provided the translation is recorded.  When $|\mathcal B|\le\eps^{-C}$,
vanishing exceptional probability is sufficient for the last term.

\begin{lemma}[Domination and conditioning on typical cells]
\label{lem:typical-conditioning}
Let $\mathsf U,\mathsf V$ be probability laws on the same finite space
with $\mathsf U\le A\mathsf V$, $A\ge1$, and let $C$ be a finite observation.
For $0<\eta<1$,
\begin{equation}
 \mathsf U_C\{c:\mathsf U_C(c)<\eta\mathsf V_C(c)\}\le\eta.
 \label{eq:conditioning-likelihood}
\end{equation}
At every other $\mathsf U_C$-positive value,
\[
 \mathsf U(\,\cdot\mid C=c)
 \le \frac A\eta\,\mathsf V(\,\cdot\mid C=c).
\]
The estimate is independent of the number of values of $C$.

In particular, let $\mathsf P$ be a finite law, let $\mathcal E$ have
mass $m=\mathsf P(\mathcal E)>0$, and put
$\mathsf U=\mathsf P(\,\cdot\mid\mathcal E)$ and
$r_C(c)=\mathsf P(\mathcal E\mid C=c)$.  Then
\begin{equation}
 \begin{gathered}
 \mathsf U_C(c)=\frac{\mathsf P_C(c)r_C(c)}m,\qquad
 \mathsf U_C\{r_C<t\}\le\frac tm\quad(t>0),\\
 \mathsf U(\,\cdot\mid C=c)
 \le r_C(c)^{-1}\mathsf P(\,\cdot\mid C=c).
 \end{gathered}
 \label{eq:retention-fibers}
\end{equation}
The last inequality is asserted on positive $\mathsf U_C$ fibers.
Thus for $0<\eta_0<1$, outside $\mathsf U_C$-mass $\eta_0$, their
domination constants are at most $(m\eta_0)^{-1}$.

Suppose further that $C$ is a function of a finite parent observation
$X$, and write $r_X(x)=\mathsf P(\mathcal E\mid X=x)$.  On every
positive $C=c$ fiber, for $0<\eta_1<1$,
\begin{equation}
 \mathsf U\{r_X(X)<\eta_1r_C(c)\mid C=c\}\le\eta_1.
 \label{eq:nested-retention-fibers}
\end{equation}
Removing these parent values does not change a conditional law given
$X$.  Let $d_-\le d_+$ be real numbers.  If a finite child $Y$
satisfies, on every $(X,Y)$ value attained in $\mathcal E$,
\[
 d_-\le-\ell^{-1}\log\mathsf P(Y\mid X)\le d_+,
 \qquad \ell=\log(1/\eps),
\]
then there are at most $\eps^{-d_+}$ retained children over each
parent, and, with $\mathsf U_c=\mathsf U(\,\cdot\mid C=c)$,
\begin{equation}
 d_- -\frac{\log(1/r_C(c))}{\ell}
 \le \Ent^{\mathsf U_c}(Y\mid X)\le d_+.
 \label{eq:typical-cell-entropy}
\end{equation}
More precisely, outside $\mathsf U_c$-mass at most $\eta_1+\eta_2$,
\begin{equation}
 d_- -\frac{\log(1/(\eta_1r_C(c)))}{\ell}
 \le-\ell^{-1}\log\mathsf U_c(Y\mid X)
 \le d_+ +\frac{\log(1/\eta_2)}{\ell},
 \qquad 0<\eta_2<1.
 \label{eq:typical-cell-surprise}
\end{equation}
These conclusions concern the retained support and sampled
conditional values.  They assert no lower mass bound for every
untrimmed child.
\end{lemma}

\begin{proof}
Summing $\mathsf U_C(c)<\eta\mathsf V_C(c)$ over the indicated
values proves \eqref{eq:conditioning-likelihood}.  On its complement,
divide $\mathsf U(\omega)\le A\mathsf V(\omega)$ by
$\mathsf U_C(c)\ge\eta\mathsf V_C(c)$.  This proves the first part.
For an event restriction the displayed identity for $\mathsf U_C$
is exact.  Summing it over $r_C<t$ proves the second assertion in
\eqref{eq:retention-fibers}; conditioning on $\mathcal E$ inside one
$C$ fiber proves the third.

Because $C$ is determined by $X$, on a positive $C=c$ fiber we have
\[
 \mathsf U_{X\mid c}(x)
 =\mathsf P_{X\mid c}(x)\frac{r_X(x)}{r_C(c)}.
\]
Summing over $r_X<\eta_1r_C(c)$ gives
\eqref{eq:nested-retention-fibers}.  The original conditional child
law on this fiber is still $\mathsf P(Y\mid X)$.
The upper surprise bound on $\mathcal E$ gives
$\mathsf P(y\mid x)\ge\eps^{d_+}$ for every retained child, so summing
these masses gives the support count.  Apply
Lemma~\ref{lem:entropy-restriction} to
$\mathsf U_c\le r_C(c)^{-1}\mathsf P(\,\cdot\mid C=c)$ for the lower
bound in \eqref{eq:typical-cell-entropy}; the support count gives its
upper bound.

For a retained parent,
\[
 \mathsf U_c(y\mid x)
 =\frac{\mathsf P(Y=y,\mathcal E\mid X=x)}{r_X(x)}
 \le\frac{\mathsf P(y\mid x)}{r_X(x)}.
\]
On the parents retained in \eqref{eq:nested-retention-fibers}, this
is at most $\eps^{d_-}/(\eta_1r_C(c))$.
On any parent fiber the support count bounds by $\eta_2$ the
$\mathsf U_c$-mass of children with
$\mathsf U_c(y\mid x)<\eta_2\eps^{d_+}$.  The two exceptional sets
give \eqref{eq:typical-cell-surprise}.
\end{proof}

\begin{lemma}[Finite splitting and resampling]
\label{lem:entropy-splitting}
For any finite tag $Q$,
\begin{align*}
 0\le\Entn(D\mid B)-\Entn(D\mid B,Q)&\le\Entn(Q),\\
 |\Mutn(T;D\mid Q)-\Mutn(T;D)|&\le\Entn(Q),\\
 \Mutn(U;V\mid Z,Q)&\le\Mutn(U;V\mid Z)+\Entn(Q).
\end{align*}
The conditional quantities on the left are averages of the quantities
on the individual tag fibers.  If $U,V$ are conditionally independent
given $Z$ under $\mathsf P$, and $\mathsf R\le K\mathsf P$, then
\[
 \Mutn^{\mathsf R}(U;V\mid Z)\le\log K.
\]
Given a joint law of $(T,C,V)$ and a kernel $\kappa(u\mid T,C)$,
one may adjoin $U$ by this kernel while preserving that entire joint
law.  The new $U$ is conditionally independent of $V$ given $(T,C)$,
and its joint law with $(T,C)$ is the prescribed kernel law.
\end{lemma}

\begin{proof}
The first difference equals $\Mutn(D;Q\mid B)$.
For the second, the chain rule gives
\[
 \Mutn(T;D\mid Q)-\Mutn(T;D)
 =\Mutn(T;Q\mid D)-\Mutn(T;Q),
\]
the difference of two numbers in $[0,\Entn(Q)]$.
Likewise
\[
 \Mutn(U;V\mid Z,Q)-\Mutn(U;V\mid Z)
 =\Mutn(U;Q\mid V,Z)-\Mutn(U;Q\mid Z)
 \le\Entn(Q).
\]
For the dominated law, decompose conditional relative entropy against
$\mathsf P_{U\mid Z}\otimes\mathsf P_{V\mid Z}$ into conditional
mutual information plus the two nonnegative marginal relative
entropies.  Its average is at most
$\KL(\mathsf R_{U,V,Z}\Vert\mathsf P_{U,V,Z})\le\log K$.
Finally the resampled joint law is exactly
$\mathsf P(t,c,v)\kappa(u\mid t,c)$; summing and conditioning prove
all its stated properties.
\end{proof}

\begin{lemma}[Finite homogenization]
\label{lem:homogenization}
Fix a law $\mathsf P$, a finite list $(X_j,Y_j)$, $1\le j\le q$,
and numbers $A_j\ge0$ such that $|\mathcal Y_j|\le\eps^{-A_j}$.
There is no cardinality assumption on $X_j$.
For $a,\sigma,\theta>0$, there are disjoint events $E_z$ whose union
has probability at least $1-q\eps^\sigma-\theta$, with the following
properties.  The number of events is at most
\[
 J=\prod_{j=1}^q\left(1+\left\lceil\frac{A_j+\sigma}{a}\right\rceil\right),
 \qquad m_z:=\mathsf P(E_z)\ge\theta/J.
\]
For each $z,j$ there is $d_{j,z}\in[0,A_j+\sigma]$ such that, on $E_z$,
\begin{equation}
 d_{j,z}\le -\ell^{-1}\log\mathsf P(Y_j\mid X_j)
 \le d_{j,z}+a.
 \label{eq:homogeneous-bin}
\end{equation}
If $\mathsf R$ is any law supported on $E_z$ satisfying
$\mathsf R\le K\mathsf P(\,\cdot\mid E_z)$, then
\begin{equation}
 d_{j,z}-\frac{\log(K/m_z)}\ell
 \le\Ent_{\mathsf R}(Y_j\mid X_j)\le d_{j,z}+a.
 \label{eq:homogeneous-entropy}
\end{equation}
More precisely, for every $0<\eta<1$, outside a set of
$\mathsf R$-probability at most $2\eta$,
\begin{equation}
 d_{j,z}-\frac{\log(K/(m_z\eta))}\ell
 \le-\ell^{-1}\log\mathsf R(Y_j\mid X_j)
 \le d_{j,z}+a+\frac{\log(1/\eta)}\ell.
 \label{eq:homogeneous-surprise}
\end{equation}
There are at most $\eps^{-(d_{j,z}+a)}$ retained $Y_j$ values over
each fixed $X_j$ value.  The support-count assertion and
\eqref{eq:homogeneous-bin} are uniform on the retained support,
before any further restriction is chosen.
\end{lemma}

\begin{proof}
For every $x$,
\[
 \sum_{y:\,\mathsf P(y\mid x)<\eps^{A_j+\sigma}}
       \mathsf P(y\mid x)\le|\mathcal Y_j|\eps^{A_j+\sigma}
       \le\eps^\sigma.
\]
Discard these tails for the $q$ observations and partition the remaining
surprise values into intervals of length $a$.  This produces at most
$J$ cells.  Discard cells of mass less than $\theta/J$; their total mass
is at most $\theta$.  This proves the first assertions.
For each retained $x,y$, \eqref{eq:homogeneous-bin} gives
$\mathsf P(y\mid x)\ge\eps^{d_{j,z}+a}$, so summing over retained $y$
gives the stated count.  It follows that the upper entropy bound in
\eqref{eq:homogeneous-entropy} holds under every law on that support.
The lower bound follows from \eqref{eq:conditional-kl}, since
$\mathsf R\le(K/m_z)\mathsf P$ and the original surprise is at least
$d_{j,z}\ell$.

For the pointwise upper bound on conditional masses, discard $x$ for
which $\mathsf R_{X_j}(x)<\eta\mathsf P_{X_j}(x)$; their total
$\mathsf R$-mass is at most $\eta$.  At all remaining $x$,
\[
 \mathsf R(y\mid x)
 \le\frac{K}{m_z\eta}\mathsf P(y\mid x)
 \le\frac{K}{m_z\eta}\eps^{d_{j,z}}.
\]
For the reverse surprise bound, on each fiber the total
$\mathsf R$-mass of values with
$\mathsf R(y\mid x)<\eta\eps^{d_{j,z}+a}$ is at most $\eta$ by the
support count.  Averaging in $x$ and combining the two exceptional
sets proves \eqref{eq:homogeneous-surprise}.
\end{proof}

Joint bins give a sharper version for increments.  Suppose $G$ is any
parent, possibly containing $T$, and on retained triples $(g,b,d)$,
\[
 \eps^{b_0+a}\le\mathsf P(B=b\mid G=g)\le\eps^{b_0-a},
 \qquad
 \eps^{d_0+a}\le\mathsf P(B=b,D=d\mid G=g)\le\eps^{d_0-a}.
\]
Dividing these inequalities gives conditional surprises for
$D\mid(G,B)$ in $[d_0-b_0-2a,d_0-b_0+2a]$ and at most
$\eps^{-(d_0-b_0+2a)}$ retained children per $(g,b)$.
Indeed, the sum of the lower joint masses cannot exceed the upper
parent mass.  Lemma~\ref{lem:entropy-restriction} therefore bounds
entropy errors under a density loss $K$ by $2a+\log K/\ell$,
independently of the size or entropy of the parent.
These are estimates for retained support and for sampled conditional
values; no assertion about every untrimmed conditional fiber is needed.

\subsection{From small information to explicit product comparisons}

\begin{lemma}[Zero information and product domination]
\label{lem:zero-information}
Let $\pi$ be a law and let $\Lambda$ be a comparison probability law
with $\pi\ll\Lambda$.  Write $L=d\pi/d\Lambda$ and
$D=\KL(\pi\Vert\Lambda)$.  For every $v>0$,
\begin{equation}
 \pi(\log L>v)\le\frac{D+1/e}{v}.
 \label{eq:likelihood-tail}
\end{equation}
Consequently, if $D=o(\ell)$, there is $v=o(\ell)$ such that
$\pi(\log L>v)=o(1)$, and the remaining subprobability is bounded
by $e^v\Lambda$.  This applies in particular to
\[
 \Lambda(a,b,c)=\pi_C(c)\pi_{A\mid C}(a\mid c)
                            \pi_{B\mid C}(b\mid c),
 \qquad D=\Mutn^\pi(A;B\mid C).
\]
Conversely, if a probability law $\rho$ satisfies
$\rho\le K\prod_{i=1}^k\mu_i$, then
\begin{equation}
 \KL\left(\rho\middle\Vert\prod_i\rho_i\right)\le\log K.
 \label{eq:product-information}
\end{equation}
\end{lemma}

\begin{proof}
Since $x\log(1/x)\le1/e$ on $0\le x\le1$,
$\mathbb E_\pi(\log L)_-\le1/e$.
Thus $\mathbb E_\pi(\log L)_+\le D+1/e$, and Markov's inequality
gives \eqref{eq:likelihood-tail}.  For example, when $\ell\to\infty$
and $D=o(\ell)$ take $v=\sqrt{(D+1)\ell}$.
For the converse, expand
\[
 \KL\left(\rho\middle\Vert\prod_i\mu_i\right)
 =\KL\left(\rho\middle\Vert\prod_i\rho_i\right)
       +\sum_i\KL(\rho_i\Vert\mu_i)\le\log K.\qedhere
\]
\end{proof}

\begin{lemma}[Simultaneous replacement by actual marginals]
\label{lem:marginal-core}
Let $\nu$ be a subprobability of mass $m>0$ on a finite space.
For $1\le l\le r$, suppose a tuple of observations
$(F_{l1},\ldots,F_{lk_l})$ satisfies
\begin{equation}
 \nu_{F_{l1},\ldots,F_{lk_l}}
       \le A_l\prod_{j=1}^{k_l}\mu_{lj},
 \label{eq:core-comparison}
\end{equation}
where every $\mu_{lj}$ is a specified probability law.
Put $q=\sum_l k_l$.  Repeated occurrences of an observation with
different comparison laws are counted separately.  For
$0<\eta<m/q$ there is a restriction $\nu'$ of $\nu$, of mass
$M\ge m-q\eta$, such that $\rho=\nu'/M$ satisfies all the bounds
\begin{equation}
 \rho_{F_{l1},\ldots,F_{lk_l}}
 \le A_l M^{k_l-1}\eta^{-k_l}
           \prod_{j=1}^{k_l}\rho_{F_{lj}}.
 \label{eq:core-actual}
\end{equation}
Support counts persist, and every marginal mass upper bound for
$\nu$ persists for $\rho$ with the factor $M^{-1}$.
In particular, taking $\eta=m/(2q)$ retains mass at least $m/2$.
\end{lemma}

\begin{proof}
Starting with $\nu$, repeatedly remove all points with
$F_{lj}=v$ whenever its current positive marginal is less than
$\eta\mu_{lj}(v)$.  Once a value is removed its marginal remains zero,
so this procedure terminates.  Charge the mass removed at such a step
to the pair $((l,j),v)$.  Each pair is charged at most once, by less
than $\eta\mu_{lj}(v)$; summing proves the loss bound $q\eta$.
On every surviving factor value,
$\mu_{lj}(v)\le\nu'_{F_{lj}}(v)/\eta
=M\rho_{F_{lj}}(v)/\eta$.
The old upper bounds hold for the restricted subprobability.  Insert
these factor inequalities into \eqref{eq:core-comparison} and divide
by $M$ to obtain \eqref{eq:core-actual}.  Restriction can only remove
support, and normalization has precisely the claimed effect on masses.
\end{proof}

The next theorem specifies the laws compared when labels are fixed.
A point may be an exact finite particle or a tuple of spatial
observations.  Additional finite tags may be included in a label.

\begin{theorem}[Finite stars and conditioning with named comparison laws]
\label{thm:finite-conditioning}
Let $\pi(p,b)$ be a probability law with
$\mu=\pi_P$, $\lambda=\pi_B$.  Let $E$ be an allowed set with
$\pi(E)\ge1-\zeta$ and suppose
\begin{equation}
 \boldsymbol1_E(p,b)\pi(p,b)\le K\mu(p)\lambda(b),\qquad K\ge1.
 \label{eq:pair-product}
\end{equation}
Choose $0<\alpha<1$ with $\zeta<1-\alpha$, retain only point fibers
with $r(p):=\pi(E\mid P=p)\ge\alpha$, and normalize the pairs in
these fibers and in $E$ to a law $q$.  Its retained mass $M_0$ satisfies
\begin{equation}
 M_0\ge1-\frac{\zeta}{1-\alpha},\qquad
 q(p,b)\le M_0^{-1}\pi(p,b),\qquad
 q(b\mid p)\le c\lambda(b),\quad c:=K/\alpha.
 \label{eq:star-kernel}
\end{equation}
For a fixed integer $n\ge1$, sample $B_1,\ldots,B_n$
independently given $P$ with kernel $q(b\mid p)$, and let
\[
 S(p,b_1,\ldots,b_n)=q_P(p)\prod_{i=1}^nq(b_i\mid p).
\]
The following are comparison probability laws on the whole star:
\begin{align}
 \Lambda_0&=q_P\otimes\lambda^{\otimes n},
 &S&\le c^n\Lambda_0,\label{eq:star-independent}\\
 \Lambda_i&=q_{P,B_i}\otimes\bigotimes_{j\ne i}\lambda_{B_j},
 &S&\le c^{n-1}\Lambda_i\quad(1\le i\le n).
 \label{eq:star-pair}
\end{align}
Here $q_{P,B_i}$ is a copy of the retained pair law $q$.
It may be replaced by the original $\pi$ in
\eqref{eq:star-pair} at the additional cost $M_0^{-1}$.

Let $R=wS/m$, where $0\le w\le1$ and $m=\mathbb E_Sw>0$.
Fix an index set $J\subset\{1,\ldots,n\}$ of \emph{fixed labels};
write $F=\{1,\ldots,n\}\setminus J$ for the \emph{free labels}.
Choose any observation
\[
 C=(B_J,Z),\qquad Z=f(P,B_J).
\]
Thus the function defining $Z$ involves the point and the fixed labels
only.  For every $0<\eta<1/(n+1)$, there is a set of $C$ values of
$R_C$-mass at least $1-(n+1)\eta$ on which, simultaneously,
\begin{align}
 R(\,\cdot\mid C)&\le\frac{c^n}{m\eta}\Lambda_0(\,\cdot\mid C),
 \label{eq:conditional-independent}\\
 R(\,\cdot\mid C)&\le\frac{c^{n-1}}{m\eta}
                 \Lambda_i(\,\cdot\mid C),\quad1\le i\le n.
 \label{eq:conditional-pair}
\end{align}
For a fixed value $C=(b_J,z)$, put $z_{b_J}(p)=f(p,b_J)$.
Ignoring the deterministic coordinates $B_J$, these comparison laws
are exactly
\begin{align}
 \Lambda_0(\,\cdot\mid b_J,z)
 &=q_P(\,\cdot\mid z_{b_J}(P)=z)
                            \otimes\bigotimes_{j\in F}\lambda_{B_j},
 \label{eq:fixed-comparison-zero}\\
 \Lambda_i(\,\cdot\mid b_J,z)
 &=q_P(\,\cdot\mid B_i=b_i,z_{b_J}(P)=z)
                            \otimes\bigotimes_{j\in F}\lambda_{B_j},
                  &&i\in J,\label{eq:fixed-comparison-pair}\\
 \Lambda_i(\,\cdot\mid b_J,z)
 &=q_{P,B_i}(\,\cdot\mid z_{b_J}(P)=z)
                       \otimes\bigotimes_{j\in F\setminus\{i\}}\lambda_{B_j},
                  &&i\in F.\label{eq:free-comparison-pair}
\end{align}
Conditionals in this display need only be defined when their
conditioning event has positive comparison mass; the inequalities
hold at every $R$-positive value under consideration.

There is also a direct comparison using the actual point marginal.
There is a set $G$ of $(P,C)$ values with $R(G)\ge1-\eta$ such that,
for every $C$ value, as an inequality of subprobabilities,
\begin{equation}
 \boldsymbol1_G(p,C)R(p,b_F\mid C)
 \le\frac{c^{|F|}}{m\eta}R_P(p\mid C)
                             \prod_{j\in F}\lambda(b_j).
 \label{eq:actual-point-free}
\end{equation}
For $0<\tau<1$, outside $C$ values of $R_C$-mass at most
$\eta/\tau$, the subprobability on the left has mass at least
$1-\tau$.  In any of these conditionals, all applicable product
comparisons, including product comparisons for coordinate tuples at
several fixed labels, may be converted \emph{simultaneously} to
comparisons using their actual retained marginals by
Lemma~\ref{lem:marginal-core}.  If the subprobability after all prior trims, including intersection
with $G$ when used, has conditional mass $u$ and the resulting list has
$N_{\mathrm{fac}}$ factor occurrences, its core has conditional mass
at least $u/2$ and each $k$-factor comparison constant $A$ becomes at
most $A(2N_{\mathrm{fac}}/u)^k$.
All exceptional-mass and domination estimates above are independent
of the number of possible values of $C$, $P$, and the labels.
\end{theorem}

\begin{proof}
The probability under $\mu$ of $r(P)<\alpha$ is at most
$\zeta/(1-\alpha)$.  More exactly, on that set,
$r\le\alpha(1-r)/(1-\alpha)$, whence
\[
 M_0=\mathbb E_\mu r-\mathbb E_\mu[r\boldsymbol1_{r<\alpha}]
 \ge1-\zeta-\frac{\alpha\zeta}{1-\alpha}
 =1-\frac{\zeta}{1-\alpha}.
\]
On retained point fibers,
$q(b\mid p)=\pi(b\mid p)\boldsymbol1_E(p,b)/r(p)
\le(K/\alpha)\lambda(b)$.
Multiplying the $n$ kernel bounds proves
\eqref{eq:star-independent}.  Leaving the $i$th pair unestimated and
bounding the other $n-1$ kernels proves \eqref{eq:star-pair}.

Apply Lemma~\ref{lem:typical-conditioning} to $\mathsf U=R$ and each of
$\mathsf V=\Lambda_0,\ldots,\Lambda_n$, with
$A=c^n/m$ or $c^{n-1}/m$.  A union bound proves
\eqref{eq:conditional-independent}--\eqref{eq:conditional-pair}.
The factorizations
\eqref{eq:fixed-comparison-zero}--\eqref{eq:free-comparison-pair}
follow by inserting $B_J=b_J$ and $z_{b_J}(P)=z$ in the displayed
product laws.  No factor depending on a free label is conditioned
on except the pair factor in \eqref{eq:free-comparison-pair}.

For \eqref{eq:actual-point-free}, independence in $S$ gives
\[
 S(p,b_F,C)=S_{P,C}(p,C)\prod_{j\in F}q(b_j\mid p).
\]
Apply \eqref{eq:conditioning-likelihood} to the marginals
$R_{P,C},S_{P,C}$ and set
$G=\{R_{P,C}\ge\eta S_{P,C}\}$.
Using $R\le S/m$ and the kernel bounds gives
\[
 \boldsymbol1_G R(p,b_F,C)
 \le\frac{c^{|F|}}{m\eta}R_{P,C}(p,C)
                                   \prod_{j\in F}\lambda(b_j).
\]
Divide by $R_C(C)$.  Since $R(G^c)\le\eta$, Markov's inequality
bounds by $\eta/\tau$ the mass of conditionals losing more than
$\tau$.  Apply Lemma~\ref{lem:marginal-core} to the subprobability in each
selected conditional, with its complete finite list of factor
comparisons and threshold $u/(2N_{\mathrm{fac}})$.  Its mass $M\le1$
makes the bound in \eqref{eq:core-actual} at most
$A(2N_{\mathrm{fac}}/u)^k$.
\end{proof}

We spell out two consequences of the theorem to make its uses precise.
First, take $Z$ trivial and $i\in J$.  For typical fixed values $b_J$,
\begin{equation}
 R_P(\,\cdot\mid B_J=b_J)
       \le\frac{c^{n-1}}{m\eta}q_P(\,\cdot\mid B_i=b_i).
 \label{eq:fixed-pair-inheritance}
\end{equation}
Thus if at label $b_i$ a coordinate tuple $\Phi_i(P)$ has a bound
$q_{\Phi_i\mid b_i}\le A_i\prod_l\mu_{il}$, its law after fixing
all of $B_J$ has bound
$c^{n-1}A_i/(m\eta)$ against exactly those factors.
This is true for every $i\in J$ simultaneously.  Applying the core
lemma to the combined list supplies actual marginals for all of them
in one retained law.  Counts on the original pair support are preserved.
If $Z$ is nontrivial, the comparison is instead the conditional law
$q_P(\,\cdot\mid B_i=b_i,z_{b_J}(P)=z)$ in
\eqref{eq:fixed-comparison-pair}; additional coordinate product
bounds must be conditioned by \eqref{eq:conditioning-likelihood},
with their factors explicitly identified.

Second, suppose $\rho(a,b)\le K\mu(a)\lambda(b)$ and $C=f(A)$.
Outside $\rho_C$-mass at most $\eta$,
\begin{equation}
 \rho(a,b\mid C=c)\le\frac K\eta
                      \mu(a\mid C=c)\lambda(b).
 \label{eq:parent-product}
\end{equation}
The measure on children is the original $\mu$ conditioned on its parent;
there is no factor $\mu_C(c)^{-1}$ in addition to the displayed
likelihood loss.  Any number of parent cells, in particular polynomially
many, is allowed.  Subsequent simultaneous marginal replacement uses
the actual child and parameter populations retained in that cell.
An observation depending on a free label does not have the factorization
used in \eqref{eq:actual-point-free}; it must first be included among
fixed labels or handled by a separately verified comparison law.

For example, if a final retained point law has
\[
 \rho_{X,U_0}\le A\rho_X\rho_{U_0},\qquad
 \rho_{X,N,L,Z}\le B\rho_X\rho_{N,L}\rho_Z,
\]
then conditioning the second inequality on $(X,N)=(x,n)$ gives
\begin{equation}
 \rho(l,z\mid x,n)
 \le B\frac{\rho_X(x)\rho_N(n)}{\rho_{X,N}(x,n)}
                     \rho(l\mid n)\rho_Z(z).
 \label{eq:full-precision-slice}
\end{equation}
The factor in the fraction exceeds $1/\eta$ on a set of
$\rho_{X,N}$-mass at most $\eta$.
This follows by summing $\rho_{X,N}(x,n)$ over the exceptional
values, so it is valid also for labels recorded at the finest precision.
The first displayed comparison remains available on the same law;
it is not replaced by a marginal from an earlier trim.

\begin{corollary}[Point slicing after a fixed-label core]
\label{cor:point-slicing}
Fix the values of some labels.  On that fiber let $\nu$ be a retained
probability law of a point $P$ and free labels $(B_j)_{j\in F}$,
with
\[
 \nu\le A\nu_P\otimes\bigotimes_{j\in F}\lambda_j.
\]
The comparison laws $\lambda_j$ may be the actual marginals obtained
by a simultaneous core on this fiber.  If $Z=f(P)$, then on every
positive $Z$ fiber,
\[
 \nu(p,b_F\mid Z=z)
 \le A\nu_P(p\mid Z=z)\prod_{j\in F}\lambda_j(b_j).
\]
Thus the comparison label populations are those \emph{before} the
point slice, and any mass upper bounds already proved for them remain
available.  After a further restriction and core of conditional mass
$M>0$, each new actual label marginal is bounded by
$A/M$ times its respective $\lambda_j$.
\end{corollary}

\begin{proof}
The comparison $Z$ marginal equals $\nu_Z$, so restricting to a
$Z$ fiber and dividing by this common mass proves the assertion.
Integrating the point and other
labels proves the marginal upper bound before the final restriction;
restriction and normalization cost at most $M^{-1}$.
\end{proof}

All these conclusions concern a fixed finite star.  They supply no
lower bound for the density of an arbitrary graph of cyclic incidences.
If $K$, $M_0^{-1}$, $m^{-1}$, $u^{-1}$, $\eta^{-1}$ and $\alpha^{-1}$ have
logarithms $o(\ell)$ and $n,N_{\mathrm{fac}}$ are fixed, every displayed loss is
$\eps^{-o(1)}$.  For a normalized retained incidence law
$\rho\le A\mu\otimes\lambda$, its support has
$(\mu\otimes\lambda)$-mass at least $A^{-1}$, simply by integrating
the domination inequality.  This is the dense-graph conclusion used
in the projection theorem.

\subsection{Countable profiles and the order of the limits}

The use of entropy at successive resolutions is related to the local
entropy averages and comparable-partition arguments of
\citet[Sections~3.3 and~4.2]{HochmanShmerkin2012} and
\citet[Section~3]{Hochman2014}.
Here we need a finite-law diagonal with explicit retention properties,
which we establish directly.

A conditional observation is \emph{homogeneous with exponent $d$}
if its realized conditional surprise, divided by $\ell$, converges
to $d$ in probability and its normalized conditional entropy also
converges to $d$.  When further restrictions will be made, we first
retain the uniform bins of Lemma~\ref{lem:homogenization}.
Convergence in probability by itself is insufficient for this purpose.
For example, let $\ell\to\infty$, give one atom mass
$1-e^{-\sqrt\ell}$, and distribute the remaining mass uniformly
among $\lceil e^\ell\rceil$ atoms.  Normalized surprise converges
to zero in probability, but restriction to the latter atoms has density
$e^{\sqrt\ell}$ and entropy tending to one in normalized units.
The support and bin requirements in Lemma~\ref{lem:homogenization}
exclude precisely this issue.

\begin{lemma}[Profiles, countable homogenization, and diagonals]
\label{lem:profile-diagonal}
The following three statements hold.
\begin{enumerate}
\item
Suppose for each $\eps\to0$ there is a countable list
$(X_j,Y_j)_{j\ge1}$, with
$|\mathcal Y_j|\le\eps^{-A_j}$ for each fixed $j$.
One can partition all but $o(1)$ of the law into
$\eps^{-o(1)}$ parts, each of mass at least $\eps^{o(1)}$,
such that the following holds.  From any choice of one retained
part at each scale, a subsequence can be extracted on which every
listed observation is homogeneous.  The same exponents are inherited
by every sequence of laws supported on those parts whose density
relative to their normalized laws is at most $\eps^{-o(1)}$.
The assertion is simultaneous for joint observations and observations
conditional on an exact particle index, when these are included in
the list.
\item
Let $v$ range over a compact rectangle of depth parameters.
Suppose observations $Y_v$, conditional on any $X$, obey, uniformly
for fixed $u,v$, the two ambiguity bounds
\[
 \log N(Y_v\mid Y_u,X),\ \log N(Y_u\mid Y_v,X)
       \le C\ell\|u-v\|_1+o(\ell),
\]
and have bounded normalized alphabet sizes at one reference depth.
Their normalized entropy profiles have subsequences converging on
all depths to a $C$-Lipschitz function, after extraction on a countable
dense set.  Homogeneous surprise exponents on that dense set extend
to the same function at every fixed depth, provided the ambiguity
bounds also hold on the retained support.
\item
Suppose a finite construction at each fixed $h>0$ has finitely many
normalized errors tending to zero as $\eps\to0$.  Given any sequence
of permissible $h_k\downarrow0$, and any countable sequence of such
requirements, a diagonal can be chosen so that the first $k$
requirements have errors at most $h_k/k$ and
$h_k\log(1/\eps_k)\to\infty$.
Thus errors may be made $o(h)$ before passing to the scale
$\eps^h$.  This assertion supplies no rate uniform in $h$ or in the
limiting profile.
\end{enumerate}
\end{lemma}

\begin{proof}
For the first assertion, apply Lemma~\ref{lem:homogenization} to
lists of length $q_k\uparrow\infty$, meshes $a_k\downarrow0$,
and, for example, $\sigma=1$, $\theta_k\downarrow0$.
At stage $k$ the number $J_k$ is finite and independent of $\eps$.
Choose the stages sufficiently slowly that
$q_k\eps\to0$ and
$\log(J_k/\theta_k)=o(\ell)$.
The discarded mass tends to zero and every retained part has the
claimed lower mass.  For each fixed $j$ the bin centers stay in the
compact interval $[0,A_j+1]$.  Successive extraction and a diagonal
make all those centers converge.  Equations
\eqref{eq:homogeneous-entropy} and
\eqref{eq:homogeneous-surprise}, with $\eta\downarrow0$
sufficiently slowly and $\log(1/\eta)=o(\ell)$, prove homogeneity.
The same estimates prove its asserted inheritance for any dominated
sequence; no exceptional set from before binning is reused.
All bounds depend on the child alphabets only.

For the second assertion, the chain rule bounds the difference of
conditional entropies by the two ambiguity logarithms.
A bounded sequence of profiles therefore has, by diagonal extraction
on a dense countable set, a limit satisfying the Lipschitz inequality
there.  Extend it continuously.  Approximation of any fixed parameter
by dense parameters and the same inequality proves convergence there.
For surprises, the identity
\[
 -\log\mathsf P(Y_v\mid X)+\log\mathsf P(Y_u\mid X)
 =-\log\mathsf P(Y_v\mid Y_u,X)
       +\log\mathsf P(Y_u\mid Y_v,X)
\]
is exact at sampled values.  For an observation with at most $N$
values on each conditional fiber, the conditional-surprise tail above
$\log N+t$ has probability at most $e^{-t}$, by summing masses
less than $e^{-t}/N$.  Apply this to both terms.  Since conditional
surprises are nonnegative, their difference in absolute value is
at most the larger of their two upper bounds outside probability
$2e^{-t}$.  Take, for instance, $t=\sqrt\ell$, then approximate
$v$ by the dense set.  This proves the assertion about surprise.

For the last assertion, fix $h_k$ first.  The defining small-parameter
convergences give a sufficiently small $\eps_k$ for the first $k$
errors to be at most $h_k/k$.  Decrease $\eps_k$ further so that
$h_k\log(1/\eps_k)\ge k$ and the scales decrease strictly.
This inductive choice proves the claim.  If a profile must first be
chosen and differentiated, choose its regular points and admissible
$h_k$ before making this final small-parameter selection.
\end{proof}

In applications with a block length $h$, ``subpower loss'' means
$\log K=o(h\ell)$, not merely $o(\ell)$.
For example the logarithm of a conditional $k$-factor core constant
in Theorem~\ref{thm:finite-conditioning} is bounded by a fixed
linear combination of
\[
 n\log(K/\alpha),\quad\log(1/M_0),\quad\log(1/m),\quad
 \log(1/\eta),\quad k\log(2N_{\mathrm{fac}}/u).
\]
These quantities must be chosen $o(h\ell)$ with $h$ fixed first.
The finite numbers of observations and replicas are also fixed at
that stage.  Lemma~\ref{lem:profile-diagonal} then permits increasing
them slowly and shrinking $h$ afterwards.  When an increment has a
much larger parent entropy, \eqref{eq:conditional-kl-bound} and the
uniform child counts give the same $o(h)$ error without multiplying
by that parent entropy.

Finally, the use of universal bounds after conditioning is a selection
principle for a specified array of conditional laws.  At each scale
fix a class $\mathcal A_\eps$ of laws and parameters such that every
sequence selected from these classes is admissible for the asserted
universal bound.  In particular, all vanishing-error hypotheses must
have common deterministic envelopes within the array: a hypothesis
$\log K=o(\ell)$ is represented there by
$\log K\le r_\eps\ell$ with one $r_\eps\to0$.
If $F_\eps(\mathsf P,\vartheta)$ has
$\limsup F_\eps\le0$ along every such selected sequence, then
\[
 \sup_{(\mathsf P,\vartheta)\in\mathcal A_\eps}
               F_\eps(\mathsf P,\vartheta)\le o(1).
\]
Otherwise choose a violating law and parameter at each of a sequence
of scales.  The same upper bound holds for any average over
conditional laws in this array, regardless of the number of
conditioning cells.  The envelopes may depend on the construction;
this argument asserts no uniform convergence over all possible
sequences satisfying an unspecified subpower condition.
When the universal statement is initially given for fixed parameters,
this argument requires stability under convergent parameters;
compactness alone does not supply that stability.  The geometric
applications below verify it by bounded box enlargements and
$O(\|\vartheta-\vartheta'\|)$ changes in depth weights.
